% ===================================================================
%  ТАБЛИЦА № 2 — Человеческое тело. — Перемена. — Игры.
%  Traduction 2026 d'apres texto/io, controlee sur texto/fr, texto/en
%  et texto/es. Consigne : voir texto/ru/10-tablica-01.tex.
% ===================================================================

%%K t02-apar-1 apar
\VUsaut{4.0mm}
\VUtitre{15.0pt}{60}{ТАБЛИЦА № 2}
\VUsaut{2.0mm}
\VUtitre{11.4pt}{0}{Человеческое тело. --- Перемена.}
\VUtitre{11.4pt}{0}{Игры.}

%%K t02-tit-0 sub
\VUtitre{13.2pt}{0}{Человеческое тело.}
\VUsaut{2.4mm}
\VUcentre{10.2pt}{0}{\textit{(Простой урок естествознания.)}}
\VUsaut{3.0mm}

%%K t02-01-1 p
1. --- Главные части человеческого тела: \VUgras{голова}, \VUgras{туловище}
и \VUgras{конечности}.

%%K t02-tit-1 sub
\VUpk{10.2pt}{40}{Голова.}
\VUsaut{3.0mm}

%%K t02-02-1 p
2. --- В голове две части: \VUgras{череп}, в котором помещается
\VUgras{мозг} и который покрыт \VUgras{волосами}, и \VUgras{лицо}.

%%K t02-03-1 p
3. --- На нашем рисунке не видно ни черепа, ни мозга; зато очень ясно
видны важные части лица.

%%K t02-04-1 p
4. --- Вот прежде всего \VUgras{лоб}, говорящий о сильном уме и о
постоянно работающей мысли. \VUgras{Виски} очень открыты. \VUgras{Глаз}
живой и широко раскрытый. Густая \VUgras{бровь} и длинные \VUgras{ресницы}
защищают его.

%%K t02-05-1 p
5. --- Вот также \VUgras{нос} и \VUgras{ноздри}. Нос служит нам для
обоняния и для дыхания. По обе стороны носа --- \VUgras{щёки}, а дальше ---
\VUgras{уши}. Нижнюю часть лица составляют \VUgras{рот} и \VUgras{подбородок}.

%%K t02-06-1 p
6. --- Их плохо видно, потому что \VUgras{усы} и \VUgras{борода} скрывают их
от нас. Но мы знаем, что во рту есть две \VUgras{губы}, \VUgras{верхняя
челюсть} и \VUgras{нижняя челюсть} с их \VUgras{зубами}, и \VUgras{язык}.
Ртом мы говорим и едим. Языком мы чувствуем вкус пищи.

%%K t02-07-1 p
7. --- Глаз, ухо, нос и рот вместе с пальцами --- органы пяти чувств:
\VUgras{зрения}, \VUgras{слуха}, \VUgras{обоняния}, \VUgras{вкуса} и
\VUgras{осязания}.

%%K t02-08-1 p
8. --- Голова, стало быть, одна из самых любопытных частей
человеческого тела.

%%K t02-tit-2 sub
\VUsaut{4.4mm}
\VUpk{10.2pt}{40}{Туловище.}
\VUsaut{3.0mm}

%%K t02-09-1 p
9. --- \VUgras{Шея} соединяет голову с туловищем. В неё входят
\VUgras{горло} и \VUgras{затылок}.

%%K t02-10-1 p
10. --- В туловище тоже помещается много необходимых органов. В
\VUgras{груди} находятся \VUgras{лёгкие}, которыми мы дышим, и
\VUgras{сердце}, средоточие жизни. Когда сердце перестаёт биться, живое
существо умирает.

%%K t02-11-1 p
11. --- \VUgras{Рёбра} держат грудь.

%%K t02-12-1 p
12. --- Другие части туловища --- \VUgras{бок}, \VUgras{спина},
\VUgras{плечи} и \VUgras{живот} (брюхо). Мы ложимся спать на бок или на
спину, а тяжести носим на плече.

%%K t02-13-1 p
13. --- В животе находятся \VUgras{желудок} и \VUgras{кишки}. Они служат
нам для переваривания пищи.

%%K t02-tit-3 sub
\VUsaut{4.4mm}
\VUpk{10.2pt}{40}{Конечности.}
\VUsaut{3.0mm}

%%K t02-14-1 p
14. --- У нас четыре \VUgras{конечности}: две \VUgras{руки} и две
\VUgras{ноги}. Руками мы берём, держим, поднимаем и собираем предметы;
ногами мы стоим, ходим и бегаем.

%%K t02-15-1 p
15. --- \VUgras{Плечо} соединяет руку с телом. Очень крепкая \VUgras{мышца},
двуглавая, двигает руку, которую \VUgras{локоть} соединяет с
\VUgras{предплечьем}. \VUgras{Запястье} образует сустав \VUgras{кисти},
оканчивающейся \VUgras{пальцами} и \VUgras{ногтями}. На кисти видна
\VUgras{вена} (и артерия), по которой течёт \VUgras{кровь}.

%%K t02-16-1 p
16. --- Части ноги: \VUgras{бедро}, \VUgras{колено} и \VUgras{подколенная
впадина}, \VUgras{икра}, \VUgras{мясо} которой покрыто \VUgras{кожей},
\VUgras{щиколотка} и \VUgras{ступня}. \VUgras{Кости} ноги очень толсты и
очень крепки. На конце ступни --- \VUgras{пальцы ног}. Другие части ---
\VUgras{подошва} и \VUgras{пятка}.

%%K t02-17-1 p
17. --- Таково почти полное описание человеческого тела.

%%K t02-17-2 p
\textit{Примечание.} --- \textit{С помощью оборота «у меня болит...
(голова и так далее)» учитель сможет пройти весь этот словарь заново со
стороны хорошего или дурного} \VUgras{телесного здоровья}.

%%K t02-tit-4 sub
\VUsaut{5.0mm}
\VUpk{10.2pt}{40}{Гимнастика.}
\VUsaut{2.4mm}
\VUcentre{10.2pt}{0}{\textit{(Рассказано одним из учеников.)}}
\VUsaut{3.0mm}

%%K t02-18-1 p
18. --- Чтобы быть гибкими, ловкими и здоровыми, надо упражнять ноги и
руки. Вот почему мы играем и занимаемся \VUgras{гимнастикой}.

%%K t02-19-1 p
19. --- На неделе оба эти занятия проходят в школьном дворе (см. таблицу
№ 1). Но по четвергам и воскресеньям мы выходим на большую лужайку, где
есть где побегать и повеселиться.

%%K t02-20-1 p
20. --- Наши учителя и родители иногда идут с нами; но упражнениями
руководит \VUgras{учитель гимнастики} \textsuperscript{(12)}.

%%K t02-21-1 p
21. --- На лужайке, слева от входа, стоят \VUgras{гимнастические снаряды}
\textsuperscript{(1)}: мы взбираемся по \VUgras{лестнице} \textsuperscript{(2)},
качаемся вниз головой на \VUgras{кольцах} \textsuperscript{(3)} и на
\VUgras{трапеции} \textsuperscript{(4)}, подтягиваемся на руках до верха
\VUgras{верёвочной лестницы} \textsuperscript{(5)} или \VUgras{каната с узлами}
\textsuperscript{(6)}.

%%K t02-22-1 p
22. --- Тем временем наши товарищи прыгают через \VUgras{деревянного коня}
\textsuperscript{(7)} или через \VUgras{брусья} \textsuperscript{(11)}. Другие
\VUgras{боксируют} \textsuperscript{(8)} или \VUgras{фехтуют} \textsuperscript{(9)} в
маске и с \VUgras{рапирой}. Третьи, наконец, поднимают \VUgras{гантели}
\textsuperscript{(10)}.

%%K t02-tit-5 sub
\VUsaut{4.6mm}
\VUpk{10.2pt}{40}{Игры.}
\VUsaut{3.0mm}

%%K t02-23-1 p
23. --- Вокруг гимнастов мальчики и девочки играют в разные \VUgras{игры}.
Девочки забавляются своим \VUgras{обручем} \textsuperscript{(14)}; они
\VUgras{прыгают через скакалку} \textsuperscript{(13)} или зовут братьев и
двоюродных сестёр на партию \VUgras{крокета} \textsuperscript{(15)}. Как они
рады отбить \VUgras{шар} противника своим \VUgras{деревянным молотком} как
раз тогда, когда он думает, что легко пройдёт в воротца!

%%K t02-24-1 p
24. --- Мальчики предпочитают игры побойчее. Они охотно играют в
\VUgras{кегли} \textsuperscript{(16)}, которые ловко сбивают \VUgras{шаром}
\textsuperscript{(17)}. Не пренебрегают они и \VUgras{волчком} \textsuperscript{(18)}
и \VUgras{бильбоке} \textsuperscript{(19)}, и даже \VUgras{лягушкой}
\textsuperscript{(20)}. Но их любимые игры --- \VUgras{шарики} \textsuperscript{(21)}
и \VUgras{лапта об стену} \textsuperscript{(22)}, потому что стена
\VUgras{оранжереи} \textsuperscript{(26)} как раз годится, чтобы отбивать от
неё лёгкий мяч.

%%K t02-25-1 p
25. --- Маленький Иван на своих \VUgras{ходулях} \textsuperscript{(24)} очень
ловок и прекрасно держит равновесие. Рядом с ним несколько товарищей
играют в \VUgras{прятки} \textsuperscript{(25)} в этой оранжерее и за
\VUgras{изгородью}. Другие предпочитают \VUgras{жмурки по руке}
\textsuperscript{(28)} или \VUgras{волан} \textsuperscript{(29)}, который перелетает
от одной \VUgras{ракетки} к другой.

%%K t02-noto-1 noto
\VUnotes{143.2mm}{%
(*) Детские игры часто носят чисто национальные имена. Мы назвали их на
идо как умели: то по самому действию, которое их отличает, то приняв
национальное имя той или другой страны, когда оно не слишком сбивает с
толку. Например: \VUgras{бочка} (немецкая и французская) --- это
\VUgras{лягушка} \textsuperscript{(20)} (английская), где играющие стараются
забросить свои круглые кружки в пасть металлической лягушки, разинутую
на дырявом ящике (бочке). \VUgras{Прыжок через плечи} \textsuperscript{(30)}
отличается от \VUgras{чехарды} только тем, что, перепрыгивая через
голову, играющие опираются руками о плечи товарища, тогда как в
«\VUgras{чехарде}» они опираются только о спину. --- Или ещё: одни
играющие нагибаются, а другие перескакивают через их спины, опираясь
руками о плечи; первые должны выдержать толчок не подавшись, вторые ---
прыгнуть не потеряв равновесия (см. саму таблицу).%
}

%%K t02-26-1 p
26. --- Посреди лужайки растут два дерева. У подножия одного из них
семь или восемь мальчиков затеяли игру в \VUgras{прыжок через плечи}
\textsuperscript{(30)} (*). Эта игра известна не во всех странах; но всякий
знает, что такое \VUgras{чехарда}, в которую дети сейчас не играют.

%%K t02-tit-6 sub
\VUsaut{5.0mm}
\VUpk{10.2pt}{40}{Игры на вольном воздухе.}
\VUsaut{3.4mm}

%%K t02-27-1 p
27. --- \VUgras{Жмурки} \textsuperscript{(31)} тоже очень весёлая игра; девочки
и мальчики по очереди надевают на глаза \VUgras{повязку} и ловят
неповоротливых, которые дают себя поймать.

%%K t02-28-1 p
28. --- Посреди лужайки --- вот совсем французская игра, хотя и
грубоватая: это \VUgras{медведь} \textsuperscript{(32)}, куда шумнее мирной игры
в \VUgras{змея} \textsuperscript{(33)} и даже английской игры в \VUgras{теннис}
\textsuperscript{(34)}.

%%K t02-29-1 p
29. --- Этот последний спорт --- отрада старших учеников. Наши учителя и
сами охотно за него берутся. Он требует большой ловкости и проворства, и
я очень его люблю. \VUgras{Качели} \textsuperscript{(35)} я оставляю девочкам.

%%K t02-30-1 p
30. --- Не по душе мне и другая английская игра, которая может выйти
опасной.

%%K t02-30-1b p
Я говорю о \VUgras{футболе} \textsuperscript{(36)}, радости моего старшего
брата. Я предпочитаю нашу старую игру в \VUgras{пленных} \textsuperscript{(38)},
столь же здоровую и куда менее грубую.

%%K t02-tit-7 sub
\VUsaut{4.6mm}
\VUpk{10.2pt}{40}{Велосипед и игры с мячом.}
\VUsaut{3.0mm}

%%K t02-31-1 p
31. --- Я и сам охотно объезжаю лужайку на \VUgras{велосипеде}
\textsuperscript{(37)}.

%%K t02-32-1 p
32. --- В будущем году я научусь \VUgras{ездить верхом} \textsuperscript{(39)}.
Как хороша лошадь, когда она красиво идёт шагом или рысью и берёт
препятствия галопом!

%%K t02-33-1 p
33. --- Летом мы учимся \VUgras{плавать} \textsuperscript{(40)}. Мы ныряем и
купаемся с наслаждением в чистой прохладной воде реки. Иногда в конце мы
пускаем бумажный \VUgras{шар} \textsuperscript{(41)}, который величаво
поднимается в воздух, как только наполнится горячим воздухом.

%%K t02-34-1 p
34. --- Есть и много других игр, но я никогда не кончил бы их
перечислять, и по этому короткому рассказу видно, что четверги на нашей
прекрасной лужайке совсем не скучны.

%%K t02-34-2 p
N.~B. --- \textit{Учителю нетрудно будет дополнить эту главу, описав
подробнее каждую из наиболее важных игр.}
\VUsaut{6.0mm}
\VUfilet{22mm}
